\documentclass[11pt,a4paper]{article}

\usepackage[T1]{fontenc}
\usepackage[utf8]{inputenc}
\usepackage[english]{babel}
\usepackage{lmodern}
\usepackage{microtype}
\usepackage{amsmath,amssymb,mathtools}
\usepackage{booktabs}
\usepackage{enumitem}
\usepackage{xcolor}
\usepackage{hyperref}
\usepackage[margin=2.45cm]{geometry}

\hypersetup{
  colorlinks=true,
  linkcolor=blue!50!black,
  citecolor=red!45!black,
  urlcolor=blue!50!black,
  pdftitle={nssDMRG: an entanglement-informed distributed DMRG method}
}

\newcommand{\ket}[1]{\lvert #1\rangle}
\newcommand{\bra}[1]{\langle #1\rvert}
\newcommand{\braket}[2]{\langle #1\vert #2\rangle}
\newcommand{\avg}[1]{\left\langle #1\right\rangle}
\newcommand{\Hil}{\mathcal H}
\newcommand{\Tr}{\operatorname{Tr}}
\newcommand{\OpVec}{\operatorname{Vec}}
\newcommand{\Id}{\mathbb 1}
\newcommand{\Order}{\mathcal O}

\title{\texttt{nssDMRG}: an Entanglement-Informed\\
Distributed DMRG Method}
\author{Method report}
\date{September 2026}

\begin{document}
\maketitle

\begin{abstract}
\texttt{nssDMRG} is a distributed-memory implementation of the density-matrix
renormalization group for one-dimensional interacting spin and fermionic
systems.  Its computational design follows directly from the bipartite and
symmetry-resolved structure of the many-body wavefunction.  At every DMRG
step, the superblock state is represented as a collection of matrices indexed
by conserved Abelian quantum numbers.  Distributing the columns of these
matrices across MPI processes makes one class of Hamiltonian contractions
strictly local, while the complementary contractions are evaluated using
communication-efficient distributed transpositions.  The full superblock
Hamiltonian need not be assembled: its action on a vector is reduced to
sector-restricted block operations that form the kernel of a Krylov
eigensolver.  This report summarizes the mathematical construction, the DMRG
workflow, and the main computational choices---sparse versus direct
application, cached versus lazy filtering, symmetry-aware communicators,
boundary conditions, and persistent storage---without entering the details
of individual source-code procedures.
\end{abstract}

\tableofcontents

\section{Overview}

The density-matrix renormalization group (DMRG) is an adaptive variational
method for quantum lattice systems.  Its efficiency in one dimension follows
from the fact that physically relevant low-energy states often have limited
bipartite entanglement.  Instead of retaining a fixed local basis as the
system grows, DMRG repeatedly selects the states carrying the dominant
Schmidt weight across a cut.

The same entanglement cut also provides a natural distributed representation
of the wavefunction.  This observation is the organizing principle of
\texttt{nssDMRG}: the computational data layout is not imposed independently
of the many-body state, but is inherited from its product-basis and symmetry
structure.  The representation serves two purposes simultaneously:
\begin{enumerate}[label=(\roman*)]
  \item it identifies the optimal retained states through the reduced density
        matrix;
  \item it decomposes the dominant Hamiltonian--vector product into local
        contractions and a small number of structured collective
        communications.
\end{enumerate}

The library supports infinite- and finite-system DMRG, Abelian quantum-number
conservation, spin and Hubbard-type fermionic models, open and periodic
boundary conditions, MPI-distributed superblock diagonalization, restart and
checkpoint facilities, and post-processing measurements based on stored
renormalization data.

\section{Lattice problem and symmetry decomposition}

Consider a finite lattice with Hilbert space
\begin{equation}
  \Hil=\bigotimes_{i=1}^{S}\Hil_i
\end{equation}
and a finite-range, or more generally $K$-local, Hamiltonian
\begin{equation}
  H=\sum_{n=1}^{K}\sum_{i_1,\ldots,i_n}
  h^{(n)}_{i_1,\ldots,i_n},
  \qquad
  h^{(n)}_{i_1,\ldots,i_n}\propto
  \prod_{m=1}^{n}O_{i_m}.
\end{equation}
An entanglement cut partitions the lattice into left and right components,
\begin{equation}
  \Hil=\Hil_L\otimes\Hil_R.
\end{equation}

We assume an Abelian conserved quantity, or a set of mutually commuting
quantities, denoted collectively by $Q$:
\begin{equation}
  [H,Q]=0.
\end{equation}
The partition is chosen so that each subsystem carries well-defined quantum
numbers.  For additive symmetries,
\begin{equation}
  q_L+q_R=Q;
\end{equation}
more generally, one may write $q_L\oplus q_R=Q$, with $\oplus$ the Abelian
composition law.  Typical examples are
\begin{equation}
  q=S^z
\end{equation}
for spin chains and
\begin{equation}
  \boldsymbol q=(N_\uparrow,N_\downarrow)
\end{equation}
for number-conserving electronic systems.

Within a fixed total sector $Q$, only compatible pairs $(q_L,q_R)$ occur.
The wavefunction decomposes as
\begin{equation}
  \ket{\Psi}
  =\sum_q\sum_{l_q=1}^{d_q^L}\sum_{r_q=1}^{d_q^R}
  \Psi^q_{r_ql_q}\,
  \ket{l_q}_L\otimes\ket{r_q}_R.
  \label{eq:psi-decomp}
\end{equation}
Equation~\eqref{eq:psi-decomp} establishes an isomorphism between a
symmetry-resolved state vector and a collection of dense matrices
\begin{equation}
  \widehat\Psi^q\in\mathbb C^{d_q^R\times d_q^L}.
\end{equation}
This matrix representation is central both to entanglement truncation and to
the parallel Hamiltonian application.

\section{DMRG renormalization cycle}

\subsection{Block enlargement and superblock construction}

At a generic step, renormalized left and right blocks are enlarged by one or
more physical sites.  For open boundary conditions, the left block grows
toward the right and the right block grows toward the left.  The enlarged
blocks form a superblock,
\begin{equation}
  \Hil_{\mathrm{SB}}=\Hil_L\otimes\Hil_R,
\end{equation}
restricted to the target total quantum-number sector.

The superblock Hamiltonian has the generic bipartite form
\begin{equation}
  H_{\mathrm{SB}}
  =H_L\otimes\Id_R
   +\Id_L\otimes H_R
   +\sum_m A_m\otimes B_m.
  \label{eq:hsb}
\end{equation}
The first two terms contain the Hamiltonians internal to each block.  The
finite set of products $A_m\otimes B_m$ contains interactions crossing the
cut: spin couplings, hopping processes, density interactions, or other
model-dependent links.

In a symmetry-adapted basis, every object in Eq.~\eqref{eq:hsb} is resolved
into blocks.  A local operator $O$ carries a quantum-number shift
\begin{equation}
  dq(O)=q_{\mathrm{out}}-q_{\mathrm{in}},
  \qquad
  O:\Hil(q_{\mathrm{in}})\longrightarrow\Hil(q_{\mathrm{out}}).
\end{equation}
For each coupling $A_m\otimes B_m$, total symmetry conservation requires
\begin{equation}
  dq(A_m)+dq(B_m)=0.
  \label{eq:dq-conservation}
\end{equation}
This metadata determines exactly which rectangular symmetry blocks can
contribute, without relying on model-specific assumptions.

\subsection{Ground-state optimization}

The lowest state of $H_{\mathrm{SB}}$ is found with an iterative Krylov-space
eigensolver.  Such algorithms require repeated evaluations of
\begin{equation}
  \ket{y}=H_{\mathrm{SB}}\ket{x},
\end{equation}
but do not intrinsically require storage of the full Hamiltonian.  Since this
matrix--vector product dominates the run time of large calculations, its
representation and parallel execution determine the practical scalability of
the method.

\subsection{Reduced density matrix and truncation}

Once the optimized state is available, the left reduced density matrix is
\begin{equation}
  \rho_L=\Tr_R\ket{\Psi}\bra{\Psi}.
\end{equation}
Using Eq.~\eqref{eq:psi-decomp}, it is block diagonal:
\begin{equation}
  \rho_L=\bigoplus_q\rho_L^q,
  \qquad
  \rho_L^q=(\widehat\Psi^q)^\dagger\widehat\Psi^q.
  \label{eq:rdm}
\end{equation}
An analogous expression holds for $\rho_R$.  Diagonalizing these matrices
gives the Schmidt probabilities $\lambda_{q\alpha}$ and the corresponding
renormalization transformations.  The retained basis is selected either by a
maximum bond dimension $\chi$ or by a discarded-weight criterion.  Its
sector-resolved dimensions obey
\begin{equation}
  \sum_q\chi_q^{L}=\chi_L,
  \qquad
  \sum_q\chi_q^{R}=\chi_R.
\end{equation}

For a retained isometry $U$, every block operator is transformed according to
\begin{equation}
  O\longmapsto U^\dagger O U.
\end{equation}
The renormalized blocks are then reused in the next growth step or in the next
finite-system sweep.  The truncation error is the total discarded density
matrix weight,
\begin{equation}
  \epsilon_{\mathrm{tr}}
  =1-\sum_{(q,\alpha)\in\mathrm{kept}}\lambda_{q\alpha}.
\end{equation}

\section{Direct Hamiltonian application}

\subsection{Vectorization identity}

Let a superblock sector be represented by the matrix $X^q$ associated with a
vector $\ket{x_q}=\OpVec(X^q)$, where columns are stacked.  For compatible
matrices $A$ and $B$,
\begin{equation}
  (A\otimes B)\OpVec(X)=\OpVec(BXA^{\mathsf T}).
  \label{eq:vec-identity}
\end{equation}
This identity replaces a very large sparse Kronecker product by two much
smaller matrix contractions.

Applied to Eq.~\eqref{eq:hsb}, it yields
\begin{align}
  (\Id_L\otimes H_R)\OpVec(X)
    &=\OpVec(H_RX),\\
  (H_L\otimes\Id_R)\OpVec(X)
    &=\OpVec(XH_L^{\mathsf T}),\\
  (A_m\otimes B_m)\OpVec(X)
    &=\OpVec(B_mXA_m^{\mathsf T}).
\end{align}
The superblock Hamiltonian is therefore applied term by term, exploiting its
tensor-product and symmetry-block structure.  The direct formulation avoids
assembling $H_{\mathrm{SB}}$, whose dimension grows as the product of the two
block dimensions.

\subsection{Sector-restricted products}

Suppose that a coupling operator connects an input left sector $k$ to an
output sector $q$.  Its filtered block is
\begin{equation}
  A_m(q,k):\Hil_L(k)\longrightarrow\Hil_L(q),
\end{equation}
and the right factor performs the complementary transition dictated by
Eq.~\eqref{eq:dq-conservation}.  The corresponding contribution is
\begin{equation}
  \OpVec\!\left[B_m(q,k)X^kA_m(q,k)^{\mathsf T}\right].
  \label{eq:coupling-product}
\end{equation}
Only symmetry-compatible pairs $(q,k)$ are evaluated.  This substantially
reduces both arithmetic and storage compared with operations in the full
product basis.

\section{Entanglement-informed MPI distribution}

\subsection{Column-wise ownership}

For each symmetry block $X^q\in\mathbb C^{d_q^R\times d_q^L}$, columns are
distributed among $P$ MPI processes.  Process $p$ owns a subset of the left
basis states and the associated column panel
\begin{equation}
  X_p^q\in\mathbb C^{d_q^R\times\Theta_{q,p}},
  \qquad
  \sum_p\Theta_{q,p}=d_q^L.
\end{equation}
This decomposition is entanglement-informed because the matrix blocks are
exactly those induced by the symmetry-preserving entanglement cut.  In DMRG,
their dimensions are the sector-resolved bond dimensions selected from the
Schmidt spectrum.

The layout has an immediate computational advantage.  Left multiplication by
a right-block operator is local:
\begin{equation}
  Y_p^q=H_R^qX_p^q.
\end{equation}
Every process possesses the complete row extent needed for each of its local
columns, and no communication is required.

\subsection{Distributed transposition}

Right multiplication by a left-block operator couples columns that reside on
different processes.  Instead of gathering the full block, the method applies
a distributed transposition, denoted by $T^*$,
\begin{equation}
  \{X_p^q\}_{p=1}^{P}
  \xrightarrow{T^*}
  \{(X^{q\,\mathsf T})_p\}_{p=1}^{P}.
\end{equation}
It is implemented through an all-to-all variable-size data exchange.  After
the transposition, the formerly nonlocal contraction becomes local; a second
$T^*$ restores the original distribution:
\begin{equation}
  X^q
  \xrightarrow{T^*}
  X^{q\,\mathsf T}
  \xrightarrow{H_L^q}
  H_L^qX^{q\,\mathsf T}
  \xrightarrow{T^*}
  X^q(H_L^q)^{\mathsf T}.
\end{equation}

\subsection{Coupling terms}

For Eq.~\eqref{eq:coupling-product}, the right factor is applied first:
\begin{equation}
  C_m(q,k)=B_m(q,k)X^k.
\end{equation}
This multiplication is local and the intermediate matrix inherits the
column-wise distribution of $X^k$.  The remaining contraction is rewritten as
\begin{equation}
  C_m(q,k)A_m(q,k)^{\mathsf T}
  =\left[A_m(q,k)C_m(q,k)^{T^*}\right]^{T^*}.
\end{equation}
It therefore requires two distributed transpositions but no global assembly
of either the wavefunction or the Kronecker operator.  The complete
Hamiltonian application consists of local sparse--dense contractions plus a
small, structured set of collective communications.

\section{Entanglement, workload, and symmetry fragmentation}

The entanglement Hamiltonian and entropy associated with the cut are
\begin{equation}
  H_E=-\log\rho_L,
  \qquad
  S_L=-\Tr(\rho_L\log\rho_L).
\end{equation}
Writing the entanglement levels as $\xi_{q\alpha}=-\log\lambda_{q\alpha}$,
the weight of sector $q$ is
\begin{equation}
  W_q=
  \frac{\sum_\alpha e^{-\xi_{q\alpha}}}
       {\sum_{k,\alpha}e^{-\xi_{k\alpha}}}.
\end{equation}
Since columns of $\widehat\Psi^q$ are distributed across processes, the same
sector structure that carries the Schmidt weight also determines the
distributed workload.  A process-level weight may be defined as
\begin{equation}
  w_p=\sum_{q\in\mathcal Q_p}W_q,
  \qquad
  \mathcal Q_p=\{q:\Theta_{q,p}>0\},
\end{equation}
with inverse participation ratio
\begin{equation}
  W_P=\sum_{p=1}^{P}w_p^2.
\end{equation}
This quantity characterizes how broadly the entanglement weight is resolved
over the process-owned subspaces.

The sector dimensions are generally nonuniform.  A few large sectors coexist
with many smaller sectors, producing \emph{symmetry fragmentation}.  This
fragmentation controls load balance and the ratio between communication and
computation.  Arithmetic is governed primarily by the volume of the matrix
blocks, while distributed transposition is controlled by their communicated
surface.  Consequently, increasing the total bond dimension improves strong
scaling until local panels become too small and latency dominates.

\subsection{Symmetry-aware subcommunicators}

A sector with dimensions $d_q^L\times d_q^R$ cannot usefully employ more
processes than its available rows or columns.  Assigning all global processes
to every small sector would create idle ranks and unnecessary collective
communication.  The implementation therefore uses a sector-dependent active
process count of the form
\begin{equation}
  P_{\mathrm{active}}(q)
  =\max\!\left[
    \min(d_q^L,P),\,
    \min(d_q^R,P)
  \right],
\end{equation}
and constructs MPI subcommunicators accordingly.  Each sector is processed
only by ranks that own a nonzero share in at least one relevant layout.  This
reduces symmetry-induced imbalance and prevents small blocks from imposing
global synchronization on the full process pool.

\section{Computational modes}

The superblock eigensolver can use three principal Hamiltonian representations.

\begin{table}[ht]
  \centering
  \small
  \begin{tabular}{@{}p{0.21\linewidth}p{0.34\linewidth}p{0.36\linewidth}@{}}
    \toprule
    Mode & Persistent objects & Main trade-off \\
    \midrule
    Sparse
      & Explicit sparse $H_{\mathrm{SB}}$
      & Simple application, but potentially prohibitive superblock memory. \\
    Direct cached
      & Sector-filtered $H_L(q)$, $H_R(q)$, $A_m(q,k)$, and $B_m(q,k)$
      & Fast repeated Krylov products at the cost of a larger persistent footprint. \\
    Direct lazy
      & Full enlarged block operators, sector state lists, and inverse maps
      & Smaller persistent memory; filtered blocks are rebuilt during each application. \\
    \bottomrule
  \end{tabular}
  \caption{High-level Hamiltonian-application strategies.}
\end{table}

The sparse mode explicitly constructs the superblock Hamiltonian.  It is
useful for smaller problems and as an independent reference path.  The direct
modes exploit Eq.~\eqref{eq:vec-identity} and are preferable when explicit
superblock storage becomes limiting.

In cached direct mode, all symmetry-filtered operator blocks needed by the
Krylov iteration are prepared once and reused.  This minimizes filtering
overhead but may consume memory comparable to an additional distributed copy
of the enlarged operator content.

In lazy direct mode, only the enlarged operators and lightweight maps between
full and sector-local indices remain resident.  A required block
\begin{equation}
  O(q,k)=O\big|_{\Hil(k)\rightarrow\Hil(q)}
\end{equation}
is constructed immediately before use and released afterwards.  The mode
trades repeated sparse filtering for a substantially smaller memory footprint,
which permits larger retained bases in memory-bound calculations.  Cached and
lazy modes share the same mathematical Hamiltonian action and MPI data layout.

\section{Spin and fermionic models}

For a spin chain, a representative nearest-neighbor Hamiltonian is
\begin{equation}
  H=J_z\sum_i S_i^zS_{i+1}^z
    +\frac{J_{xy}}{2}\sum_i
     \left(S_i^+S_{i+1}^-+S_i^-S_{i+1}^+\right).
\end{equation}
The connecting terms naturally take the form $A_m\otimes B_m$, and their
quantum-number shifts cancel pairwise.

For a Hubbard chain,
\begin{equation}
  H=-\sum_{i,\sigma}t_{i\sigma}
    \left(c_{i\sigma}^\dagger c_{i+1,\sigma}+\mathrm{H.c.}\right)
    +U\sum_i n_{i\uparrow}n_{i\downarrow},
\end{equation}
fermionic anticommutation must be respected across tensor-product boundaries.
With the convention used by the code, a boundary hopping factor on the left
is represented using the parity operator $P=(-1)^N$, for example
\begin{equation}
  (c_{L\sigma}^\dagger P_L)\otimes c_{R\sigma}.
\end{equation}
The parity factor is prepared consistently in both cached and lazy direct
modes.  Its presence changes the algebraic operator content but not the
distributed contraction pattern.

\section{Open and periodic boundary conditions}

With open boundaries, each block grows toward the central cut and the
superblock contains a single set of links connecting the enlarged blocks.
Periodic boundaries add couplings closing the two outer ends.  A naive growth
scheme would leave one boundary operator buried under many renormalization
steps, degrading its representation.

The periodic strategy therefore alternates block growth on the two sides.
Each block carries operators associated with both link directions, and the
superblock is connected at both ends.  This preserves a more balanced
renormalization history for the two periodic links.  The physical ordering of
sites is tracked separately from their growth order, a distinction that is
also required when reconstructing nonlocal observables.

\section{Memory management, restart, and measurements}

Large DMRG calculations must control persistent storage as carefully as the
Hamiltonian application.  The implementation separates several classes of
data:
\begin{itemize}
  \item renormalized block objects required to continue or restart a run;
  \item rotation matrices $U_n$ required by finite sweeps and by
        post-processing operators;
  \item stored superblock states and sector metadata required for later
        measurements;
  \item temporary sector-filtered operators used by direct products.
\end{itemize}

Block checkpoints may be written at every step or restricted to the latest
state.  Rotation matrices may remain cached in memory for speed or be stored
on disk to reduce resident memory.  Likewise, a measurement state can be
saved so that local observables, entanglement quantities, and static
correlations can be evaluated without repeating the expensive ground-state
optimization.

The use of explicit quantum-number shifts for operators extends the same
sector logic from the Hamiltonian to measurements.  A correlation
\begin{equation}
  \avg{O_A(i)O_B(j)}
\end{equation}
is allowed within a fixed total sector only if
\begin{equation}
  dq(O_A)+dq(O_B)=0.
\end{equation}
Operators residing on opposite sides of the cut are contracted with the same
rectangular, symmetry-resolved tensor-product machinery used by the direct
Hamiltonian application.  Thus the measurement layer reuses the central
mathematical abstraction rather than constructing full superblock operators.

\section{End-to-end workflow}

A complete calculation can be summarized as follows:
\begin{enumerate}[label=\arabic*.]
  \item Define the local Hilbert space, local Hamiltonian, boundary operators,
        conserved Abelian quantum numbers, and operator shifts.
  \item Initialize left and right blocks and choose the target total sector.
  \item Enlarge both blocks and construct the symmetry-compatible superblock
        basis.
  \item Select sparse, cached-direct, or lazy-direct Hamiltonian application.
  \item Distribute every wavefunction sector matrix by columns and build
        sector-specific MPI communicators.
  \item Solve the superblock eigenproblem through repeated distributed
        Hamiltonian--vector products.
  \item Form the symmetry-block reduced density matrices, diagonalize them,
        and select the retained Schmidt states.
  \item Renormalize the blocks and all required operators.
  \item Repeat the growth process for infinite DMRG, or move the cut through
        the lattice during finite sweeps.
  \item Store checkpoints, rotation matrices, and measurement states according
        to the selected memory policy.
  \item Evaluate observables and correlations in post-processing when desired.
\end{enumerate}

\section{Scaling picture and practical regime}

For sufficiently large sector matrices, local contractions dominate the
cost, while the distributed transpositions introduce a subleading
communication contribution.  The ideal strong-scaling behavior is
approximately
\begin{equation}
  T(P)\sim\frac{T(1)}{P},
\end{equation}
until one of three effects becomes important:
\begin{enumerate}[label=(\roman*)]
  \item symmetry fragmentation leaves too little work in small sectors;
  \item the local matrix panels approach cache- or latency-dominated sizes;
  \item filtering and collective-communication overhead becomes comparable
        with arithmetic.
\end{enumerate}
Increasing the bond dimension generally improves parallel efficiency because
it increases the local volume-to-communication-surface ratio.  The relevant
limit is consequently not only the formal complexity of DMRG, but also the
entanglement- and symmetry-controlled distribution of the retained basis.

The approach is not conceptually restricted to DMRG.  It requires a
symmetry-preserving bipartition, a wavefunction represented in the associated
product basis, and a Hamiltonian expressible as a manageable sum of products
across the cut.  DMRG provides a particularly natural realization because its
Schmidt decomposition already constructs and truncates precisely this basis.

\section{Summary}

The central design principle of \texttt{nssDMRG} is that entanglement and
symmetry organize both the physical approximation and the distributed
computation.  The entanglement cut produces a set of symmetry-resolved
wavefunction matrices.  These matrices define the retained DMRG basis, the
MPI ownership layout, the allowed operator transitions, and the communication
pattern of the Hamiltonian application.

The superblock Hamiltonian is reduced to local contractions of the form
\begin{equation}
  BXA^{\mathsf T},
\end{equation}
with distributed transpositions used only when the contraction direction is
incompatible with the current data layout.  Symmetry-sector filtering removes
forbidden work, while sector-specific communicators prevent fragmented blocks
from involving idle processes.  Cached and lazy modes expose an explicit
time--memory trade-off without changing the underlying algorithm.  Together
with symmetry-aware truncation, periodic-growth strategies, restart support,
and post-processing measurements, this yields a scalable DMRG framework in
which computational organization follows directly from the quantum structure
of the state.

\begin{thebibliography}{9}

\bibitem{White1992}
S.~R. White,
``Density matrix formulation for quantum renormalization groups,''
\emph{Physical Review Letters} \textbf{69}, 2863 (1992).

\bibitem{White1993}
S.~R. White,
``Density-matrix algorithms for quantum renormalization groups,''
\emph{Physical Review B} \textbf{48}, 10345 (1993).

\bibitem{Schollwock2011}
U. Schollw\"ock,
``The density-matrix renormalization group in the age of matrix product
states,''
\emph{Annals of Physics} \textbf{326}, 96 (2011).

\bibitem{Amaricci2022}
A. Amaricci \emph{et al.},
``EDIpack: A parallel exact diagonalization package for quantum impurity
problems,''
\emph{Computer Physics Communications} \textbf{273}, 108261 (2022).

\bibitem{MethodManuscript}
A. Amaricci,
``Entanglement-informed distributed wavefunction approach to scalable quantum
many-body systems,'' manuscript (2026).

\end{thebibliography}

\end{document}
